\begin{recipe}
[% 
    preparationtime = {\unit[35]{min}},
    portion = {\portion{2}},
    source = {\protect\recipesource{https://www.instagram.com/p/DObznzYCQpc/}{Instagram: wemocooks}},
    calory = {1025kcal | 28g Protein | 110g KH | 55g Fett},
    price = {3,70€},
    rating = {3},
]
{Pumpkin Gnocchi mit cremiger Kürbis-Salbei-Sauce}

\graph{% pictures
   big=pic/hauptgericht/09_pumpkin-gnocchi
}

\introduction{%
Herbstlicher geht es kaum: kleine Kürbis-Gnocchi aus Süßkartoffel, außen knusprig angebraten und serviert in einer cremigen Kürbis-Salbei-Sauce mit Parmesan.
}

\ingredients{
    \textbf{Gnocchi} & \\
    280 g & Süßkartoffel\index{Gemüse!Süßkartoffel}\\
    15 g & Kürbispüree\index{Gemüse!Kürbis}\\
    90 g & Weizenmehl\index{Getreide!Weizenmehl}\\
    1 TL & Salz\\
    2--3 EL & Olivenöl\\
    2 EL & Kürbiskerne\index{Nüsse \& Samen!Kürbiskern}\\
    \\
    \textbf{Sauce} & \\
    30 g & Butter\index{Milchprodukte \& Alternativen!Butter}\\
    1 & Zwiebel\index{Gemüse!Zwiebel}\\
    3 Zehen & Knoblauch\index{Gemüse!Knoblauch}\\
    3 EL & Kürbispüree\\
    60 ml & Sahne\index{Milchprodukte \& Alternativen!Sahne}\\
    120 ml & Vollmilch\index{Milchprodukte \& Alternativen!Milch}\\
    120 ml & Wasser\\
    30 g & Parmesan\index{Milchprodukte \& Alternativen!Parmesan}\\
    & 1 TL Zwiebelpulver\\
    & 1 TL Knoblauchpulver\\
    & 1 TL Salbeipulver\\
    & \nicefrac{1}{4} TL Muskat\\
    & \nicefrac{1}{2} TL Schwarzer Pfeffer
}

\preparation{%
    \step%
    Süßkartoffel mit einer Gabel einstechen und in der Mikrowelle weich garen. Halbieren und das Fruchtfleisch auslösen.
    
    \step%
    Mit Kürbispüree, Mehl und Salz zu einem Teig verkneten. Nur so lange kneten, bis alles gerade verbunden ist.
    
    \step%
    Kleine Kugeln formen und mit einem Messer rundherum leicht einschneiden, sodass sie wie kleine Kürbisse aussehen.
    
    \step%
    In leicht siedendem Wasser kochen, bis sie an die Oberfläche steigen. Abgießen.
    
    \step%
    In einer Pfanne mit Olivenöl von unten knusprig anbraten. Kürbiskerne kurz mitrösten.
    
    \step%
    Für die Sauce Butter schmelzen und die gewürfelte Zwiebel glasig dünsten. Knoblauch zugeben.
    
    \step%
    Kürbispüree, Milch, Sahne, Wasser und Gewürze einrühren und einkochen lassen, bis die Sauce leicht eindickt.
    
    \step%
    Parmesan unterrühren, bis er geschmolzen ist. Mit den Gnocchi servieren.
}

% \hint{
%     Der Teig sollte nicht zu stark geknetet werden, sonst werden die Gnocchi zäh.
% }

\end{recipe}